\documentclass[11pt,a4paper]{article}
\usepackage[margin=1in]{geometry}
\usepackage{graphicx}
\usepackage{booktabs}
\usepackage{float}
\usepackage{hyperref}
\usepackage{array}
\usepackage{caption}

\title{TableTalk Technical Test Report\\\large Narrative Voice Processing and Classification}
\author{Bowen Zhang}
\date{\today}

\begin{document}
\maketitle

\section{Overview}
This report summarizes the full pipeline run on a RAVDESS subset (200 speech recordings):
(1) audio preprocessing and feature extraction,
(2) narrative tone classification,
(3) AI-based transcription,
and (4) narrative audio retrieval simulation.

\section{Dataset and Split}
Data source: RAVDESS Emotional Speech Dataset (speech subset).
A balanced subset of 200 clips was used (8 emotion classes, 25 clips per class), consistent with the assignment recommendation of using a subset (50--200 recordings).

\begin{table}[H]
\centering
\caption{Dataset split summary}
\begin{tabular}{lccc}
\toprule
Split & Train & Validation & Test \\
\midrule
Clip count & 140 & 20 & 40 \\
\bottomrule
\end{tabular}
\end{table}

\begin{table}[H]
\centering
\caption{Label distribution (balanced subset)}
\begin{tabular}{lc}
\toprule
Label & Count \\
\midrule
angry & 25 \\
calm & 25 \\
disgust & 25 \\
fearful & 25 \\
happy & 25 \\
neutral & 25 \\
sad & 25 \\
surprised & 25 \\
\bottomrule
\end{tabular}
\end{table}

\begin{figure}[H]
\centering
\includegraphics[width=0.72\textwidth]{outputs/metrics/figures/split_distribution.png}
\caption{Split distribution}
\end{figure}

\begin{figure}[H]
\centering
\includegraphics[width=0.72\textwidth]{outputs/metrics/figures/label_distribution.png}
\caption{Label distribution}
\end{figure}

\begin{figure}[H]
\centering
\includegraphics[width=0.72\textwidth]{outputs/metrics/figures/duration_kde.png}
\caption{Duration KDE by label}
\end{figure}

\section{Task 1: Audio Processing and Feature Engineering}
Pipeline summary:
\begin{itemize}
  \item Resampling to 16 kHz, mono conversion.
  \item Peak normalization.
  \item Feature extraction: duration, pitch statistics, spectral centroid statistics, energy statistics, and MFCC statistics.
  \item Structured output for modeling: one row per clip in \texttt{data/features/features.csv}.
\end{itemize}

\section{Task 2: Narrative Tone Classification}
Candidate models were Logistic Regression and Random Forest.
Model selection on validation split:

\begin{table}[H]
\centering
\caption{Validation model comparison (macro-F1)}
\begin{tabular}{lc}
\toprule
Model & Validation macro-F1 \\
\midrule
Logistic Regression & 0.3589 \\
Random Forest & 0.2524 \\
\bottomrule
\end{tabular}
\end{table}

Best model: \textbf{Logistic Regression}.

\begin{table}[H]
\centering
\caption{Test metrics (best model)}
\begin{tabular}{lc}
\toprule
Metric & Value \\
\midrule
Accuracy & 0.3500 \\
Macro-F1 & 0.3427 \\
Weighted-F1 & 0.3427 \\
\bottomrule
\end{tabular}
\end{table}

\begin{table}[H]
\centering
\caption{Per-class test results}
\begin{tabular}{lccc}
\toprule
Class & Precision & Recall & F1 \\
\midrule
angry & 0.5000 & 0.2000 & 0.2857 \\
calm & 0.5000 & 0.4000 & 0.4444 \\
disgust & 0.4000 & 0.4000 & 0.4000 \\
fearful & 0.4286 & 0.6000 & 0.5000 \\
happy & 0.0000 & 0.0000 & 0.0000 \\
neutral & 0.5000 & 0.4000 & 0.4444 \\
sad & 0.1429 & 0.2000 & 0.1667 \\
surprised & 0.4286 & 0.6000 & 0.5000 \\
\bottomrule
\end{tabular}
\end{table}

\begin{figure}[H]
\centering
\includegraphics[width=0.68\textwidth]{outputs/metrics/figures/confusion_matrix.png}
\caption{Confusion matrix on test split}
\end{figure}

\paragraph{Short discussion.}
The 8-class setup is non-trivial with handcrafted features and a small subset size.
Current results are above random baseline ($1/8=12.5\%$) but still modest, indicating room for stronger representations (e.g., pretrained speech embeddings) and further tuning.

\section{Task 3: AI-Based Transcription}
ASR model: Whisper (\texttt{tiny}).
A total of 200 clips were transcribed and evaluated against RAVDESS canonical statements.

\begin{table}[H]
\centering
\caption{ASR evaluation}
\begin{tabular}{lc}
\toprule
Metric & Value \\
\midrule
Pairs evaluated & 200 \\
WER & 0.0658 \\
CER & 0.0359 \\
\bottomrule
\end{tabular}
\end{table}

\paragraph{Short discussion.}
The transcription pipeline is stable on this subset (200/200 processed).
The observed WER/CER indicates generally reliable recognition for the limited phrase set in RAVDESS.

\section{Task 4: Narrative Audio Retrieval (TableTalk Simulation)}
A rule-based retrieval prototype was implemented using label, duration, acoustic features, and transcript text.

\begin{table}[H]
\centering
\caption{Example retrieval queries and returned results}
\begin{tabular}{p{0.62\textwidth}c}
\toprule
Query & Hits \\
\midrule
\texttt{"calm narration >4 seconds"} & 9 \\
\texttt{"calm narration >4 seconds"} (saved demo file) & 10 \\
\texttt{"high-energy speech"} & 3 \\
\bottomrule
\end{tabular}
\end{table}

Result files:
\begin{itemize}
  \item \texttt{outputs/retrieval/results.csv}
  \item \texttt{outputs/retrieval/results\_calm.csv}
  \item \texttt{outputs/retrieval/results\_high\_energy.csv}
\end{itemize}

\paragraph{Short discussion.}
The prototype supports interpretable filtering and quick experimentation.
For broader narrative search, the next step is to add learned embeddings and a weighted ranking function.

\section{Reproducibility and Artifacts}
Key artifacts generated in this run:
\begin{itemize}
  \item Classification metrics: \texttt{outputs/metrics/test\_metrics.json}
  \item Classification report: \texttt{outputs/metrics/test\_classification\_report.json}
  \item ASR metrics: \texttt{outputs/metrics/asr\_metrics.json}
  \item Figures: \texttt{outputs/metrics/figures/*.png}
\end{itemize}

\end{document}
